\section{Evaluation}
\label{sec:evaluation}

\subsection{RQ3: Latency, Storage, and Resource Scaling}

RQ3 tests the source-derived cost model rather than assuming that worker count
alone explains scaling.  The protocol independently varies worker/publisher
count, manifest depth, live-session count and lease age, private-upper entry
count and fan-out, changed and validated paths, published logical bytes,
conflict class, and eligible-merge bytes, lines, similarity, and edit distance.
Storage cases include small edits to large files, many small files, and
incompressible controls.  All factors, warmups, repeats, and stopping rules
must be fixed before the final source freeze.

For each stage, we report latency distributions, CPU, RSS/PSS, I/O, failures,
and exact phase boundaries.  Publication additionally reports writer-lock wait
and hold time and decomposes planning, resolution/merge, hashing, copy,
synchronization, and manifest transition.  Storage reports logical, allocated,
shared, and exclusive bytes for upper, work, staging, published layers, and
lease-retained history.  Squash reports plan, build, commit, and garbage
collection separately.  No cost-model term becomes a paper result until it is
reproduced from frozen raw artifacts and analysis code.

\emph{The remaining RQ1--RQ5 methodology and all numerical results are pending
protocol and evidence lock.}
